%%%%%%%%%%%%%%%%%%%%%%%%%%%%% Define Article %%%%%%%%%%%%%%%%%%%%%%%%%%%%%%%%%%
\documentclass{article}
%%%%%%%%%%%%%%%%%%%%%%%%%%%%%%%%%%%%%%%%%%%%%%%%%%%%%%%%%%%%%%%%%%%%%%%%%%%%%%%

%%%%%%%%%%%%%%%%%%%%%%%%%%%%% Using Packages %%%%%%%%%%%%%%%%%%%%%%%%%%%%%%%%%%
\usepackage{geometry}
\usepackage{graphicx}
\usepackage{amssymb}
\usepackage{amsmath}
\usepackage{amsthm}
\usepackage{empheq}
\usepackage{mdframed}
\usepackage{booktabs}
\usepackage{lipsum}
\usepackage{graphicx}
\usepackage{color}
\usepackage{psfrag}
\usepackage{pgfplots}
\usepackage{bm}
%%%%%%%%%%%%%%%%%%%%%%%%%%%%%%%%%%%%%%%%%%%%%%%%%%%%%%%%%%%%%%%%%%%%%%%%%%%%%%%

% Other Settings

%%%%%%%%%%%%%%%%%%%%%%%%%% Page Setting %%%%%%%%%%%%%%%%%%%%%%%%%%%%%%%%%%%%%%%
\geometry{a4paper}
\graphicspath{ {Problems/} }

%%%%%%%%%%%%%%%%%%%%%%%%%% Define some useful colors %%%%%%%%%%%%%%%%%%%%%%%%%%
\definecolor{ocre}{RGB}{243,102,25}
\definecolor{mygray}{RGB}{243,243,244}
\definecolor{deepGreen}{RGB}{26,111,0}
\definecolor{shallowGreen}{RGB}{235,255,255}
\definecolor{deepBlue}{RGB}{61,124,222}
\definecolor{shallowBlue}{RGB}{235,249,255}
%%%%%%%%%%%%%%%%%%%%%%%%%%%%%%%%%%%%%%%%%%%%%%%%%%%%%%%%%%%%%%%%%%%%%%%%%%%%%%%

%%%%%%%%%%%%%%%%%%%%%%%%%% Define an orangebox command %%%%%%%%%%%%%%%%%%%%%%%%
\newcommand\orangebox[1]{\fcolorbox{ocre}{mygray}{\hspace{1em}#1\hspace{1em}}}
%%%%%%%%%%%%%%%%%%%%%%%%%%%%%%%%%%%%%%%%%%%%%%%%%%%%%%%%%%%%%%%%%%%%%%%%%%%%%%%

%%%%%%%%%%%%%%%%%%%%%%%%%%%% English Environments %%%%%%%%%%%%%%%%%%%%%%%%%%%%%
\newtheoremstyle{mytheoremstyle}{3pt}{3pt}{\normalfont}{0cm}{\rmfamily\bfseries}{}{1em}{{\color{black}\thmname{#1}~\thmnumber{#2}}\thmnote{\,--\,#3}}
\newtheoremstyle{myproblemstyle}{3pt}{3pt}{\normalfont}{0cm}{\rmfamily\bfseries}{}{1em}{{\color{black}\thmname{#1}~\thmnumber{#2}}\thmnote{\,--\,#3}}
\theoremstyle{mytheoremstyle}
\newmdtheoremenv[linewidth=1pt,backgroundcolor=shallowGreen,linecolor=deepGreen,leftmargin=0pt,innerleftmargin=20pt,innerrightmargin=20pt,]{theorem}{Theorem}[section]
\theoremstyle{mytheoremstyle}
\newmdtheoremenv[linewidth=1pt,backgroundcolor=shallowBlue,linecolor=deepBlue,leftmargin=0pt,innerleftmargin=20pt,innerrightmargin=20pt,]{definition}{Definition}[section]
\theoremstyle{myproblemstyle}
\newmdtheoremenv[linecolor=black,leftmargin=0pt,innerleftmargin=10pt,innerrightmargin=10pt,]{problem}{Problem}[section]
%%%%%%%%%%%%%%%%%%%%%%%%%%%%%%%%%%%%%%%%%%%%%%%%%%%%%%%%%%%%%%%%%%%%%%%%%%%%%%%

%%%%%%%%%%%%%%%%%%%%%%%%%%%%%%% Plotting Settings %%%%%%%%%%%%%%%%%%%%%%%%%%%%%
\usepgfplotslibrary{colorbrewer}
\pgfplotsset{width=8cm,compat=1.9}
%%%%%%%%%%%%%%%%%%%%%%%%%%%%%%%%%%%%%%%%%%%%%%%%%%%%%%%%%%%%%%%%%%%%%%%%%%%%%%%

%%%%%%%%%%%%%%%%%%%%%%%%%%%%%%% Title & Author %%%%%%%%%%%%%%%%%%%%%%%%%%%%%%%%
\title{MAM1019H Exam Questions}
\author{Alexander Cristaudo}
%%%%%%%%%%%%%%%%%%%%%%%%%%%%%%%%%%%%%%%%%%%%%%%%%%%%%%%%%%%%%%%%%%%%%%%%%%%%%%%

\begin{document}
    \maketitle
    \newpage
 
    \includegraphics[scale = 0.55]{Screenshot 2022-11-07 at 21.07.38.png}
    \\ \\ We can also prove injectivity with: (Note when we say $f(a)=f(b)$, we must mention $a,b\in D_f$) \\
    Suppose $X \ne Y$, for any $X,Y\in \mathcal{P}(\mathbb{N})$. Then, WLOG, we may assume $\exists x\in X $ such that $x \notin Y$. Then $f(x) \in g(X)$. Now if $f(x) \in g(Y)$, then $\exists y\in Y$ such that $f(y)=f(x)$. But $f(y)=f(x) \Rightarrow x=y$ and $x\notin y$ so we must have that $f(x) \notin g(Y) \Rightarrow g(X) \ne g(Y)$ \\ \linebreak
    \includegraphics[scale = 0.55]{Screenshot 2022-11-07 at 22.37.44.png};
    \linebreak
    \includegraphics[scale = 0.55]{Screenshot 2022-11-07 at 21.43.20.png} \\ \linebreak
    \includegraphics[scale = 0.55]{Screenshot 2022-11-07 at 22.05.09.png} \\
    \includegraphics[scale = 0.55]{Screenshot 2022-11-07 at 22.10.24.png} \\
    \includegraphics[scale = 0.55]{Screenshot 2022-11-07 at 22.15.38.png} \\
    Another possible way is: $f:\mathcal{P}(\mathbb{N}) \to X$ defined by $f(X)=(x_1, x_2, x_3, ...)$ where
    \[
    x_i = 
    \begin{cases}
        1 & \text{if } i\in A \\
        0 & \text{otherwise}
    \end{cases}
    \] \\ \linebreak
    When using CSB, state that $|A| \le |B|$ and $|B| \le |A| $ so $|A|=|B|$ rather than mentioning injectives
    \\ \includegraphics*[scale = 1.2]{Screenshot 2022-11-07 at 22.44.14.png} \\
    \includegraphics*[scale = 1.2]{Screenshot 2022-11-07 at 22.57.07.png}
    \\ For (3), we define $f:\mathcal{F} \to \mathcal{P}(\mathbb{N})$ by $f(g) = \{n\in \mathbb{N}:g(n)=1\}$ and show this is bijective
    \\ For (4), $f:\mathcal{F} \to \mathcal{P}(\mathbb{R})$ then $\mathbb{R} < \mathcal{P}(\mathbb{R})$ \\ We can also provide an injection $\mathbb{R} \to \mathcal{F}$, then assume a surjection exists. Then consider 
    \[
        h(x) = 
        \begin{cases}
            1 & \text{if } g(x)(x) = 0 \\
            0 & \text{otherwise}
        \end{cases}
    \]
    \\ Here, $g(x)$ returns some function and we take that function at $x$. Every $x\in \mathbb{R}$ differs to this function so not surjective
    \\ \includegraphics*[]{Screenshot 2022-11-07 at 23.09.17.png} \\ 
    Suppose $g'(a) = g'(b)$. Then $x_a = x_b \Rightarrow x_a\in g^{-1}({a})$ and $x_a \in g^{-1}(b)$. But $a$ can only map to one element, so $a=b$ \\
    \includegraphics*[scale = 0.4]{Screenshot 2022-11-07 at 23.17.00.png}
    \\ \includegraphics*[scale = 0.5]{Screenshot 2022-11-07 at 23.17.41.png}
    \\ \includegraphics*[]{Screenshot 2022-11-07 at 23.34.16.png} \\ When checking if something is well ordered, first check linearly ordered. Then try provide order isomorphism to $\mathbb{N}$. We have the implication if $(P, \le)$ is well ordered, and $P \cong Q$, then $(Q, \preceq)$ is well ordered. So if $(Q, \preceq)$ is not well ordered, there is no isomorphism. If we cannot provide isomorphism, it can still be well ordered. Try the subset method then. With isomorphic method, dont need to show linearly ordered but do for using the definition
        \includegraphics*[scale = 0.58]{Screenshot 2022-11-08 at 15.24.43.png} \\ 
        \linebreak
        Unless otherwise stated, can use $s(n)=n+1=1+n, 1\cdot n=n=n\cdot 1...$ all have been proven already + associativity, commutativity, distributivity of addition \\
        \includegraphics*[scale = 0.6]{Screenshot 2022-11-08 at 20.15.33.png} \\ If unsure, include a proper associative proof
        \\ \includegraphics*[scale = 0.6]{Screenshot 2022-11-08 at 20.20.51.png}
        \\ \includegraphics*[scale = 0.6]{Screenshot 2022-11-08 at 20.29.53.png}
        \\ \includegraphics*[scale = 1.2]{Screenshot 2022-11-08 at 20.37.06.png}
\end{document}